\documentclass{sprawozdanie-agh}

\usepackage[utf8]{inputenc}
\usepackage{matlab-prettifier}
\usepackage{siunitx}
\usepackage{booktabs,caption,ragged2e}
\usepackage[shortlabels]{enumitem}


\makeatletter

% ========= Początek Dokumentu ========= 

\begin{document}
% ========= Strona tytułowa ========= 
\przedmiot{Teoria Sterowania II: Ćwiczenia Laboratoryjne}
\tytul{Wytyczne dla Ćwiczenia 4}
\podtytul{Bezpośrednia Metoda Lapunowa}
\kierunek{Automatyke i~Robotyka}
\autor{Dariusz Cieślar, Maciej Różewicz}
\data{Kraków, 3 marca 2026}

\stronatytulowa{}

% ========= Treść Właściwa =========
\section{Cel ćwiczenia}
Celem ćwiczenia jest zapoznanie się z~metodami badania stabilności układów nieliniowych.
W czasie zajęć przedstawione i~przećwiczone zostanie bezpośrednia metoda
Lapunowa - bardzo
uniwersalna metoda, pozwalająca zbadać stabilność szerokiej gamy układów
(zarówno stacjonarnych, jak i~niestacjonarnych).
Ponadto po sprawdzeniu stabilności punktów równowagi układu badany i~estymowany będzie 
obszar przyciągania asymptotycznego za pomocą twierdzenia La Salle'a.
Dodatkowo sprawdzana będzie stabilność nie tylko punktów równowagi, ale i~cykli granicznych.

%%%%%%%%%%%%%%%%%%%%%%%%%%%%%%%%%%%%%%%%%%%%%%%%%%%%%%%%%%%% 
\section{Przebieg ćwiczenia}  
    %=======================================================
    \subsection{Bezpośrednia metoda Lapunowa i~twierdzenie La Salle'a}
    Dla poniższych układów (\ref{eq:zadania_przyklad_1}), (\ref{eq:zadania_przyklad_2}), (\ref{eq:zadania_przyklad_3}), (\ref{eq:zadania_przyklad_4}):
    \begin{itemize}
        \item zamodelować układ równań w~Simulinku (lub użyć skryptu do jego rozwiązania),
        \item wyznaczyć portrety fazowe,
        \item znaleźć punkty równowagi,
        \item zbadać stabilność punktów równowagi, stosując podane funkcje Lapunowa,
        \item dla każdej z~podanych funkcji Lapunowa znaleźć analitycznie estymatę obszaru atrakcji,
        \item symulacyjnie określić obszar atrakcji,
        \item porównać obszary odnalezione analitycznie i~symulacyjnie - przedstawić obydwa na tle portretu fazowego,
        \item porównać estymaty obszaru atrakcji uzyskane z~użyciem różnych funkcji Lapunowa - przedstawić obydwa na jednym wykresie na tle portretu fazowego.
    \end{itemize}
        %---------------------------------------------------
        \subsubsection{System 1}
        Rozważamy układ:
        \begin{equation}
            \label{eq:zadania_przyklad_1}
            \left\{
            \begin{array}{l}
                 \dot{x}_{1} = -x_{1} + 2x_{1}^{2}x_{2}\\
                 \dot{x}_{2} = -x_{2}
            \end{array}
            \right\}.
        \end{equation}
        Dla tego układu rozpatrujemy dwóch kandydatów na funkcję Lapunowa:
        \begin{enumerate}[label=(\alph*)]
            \item $V_{1}^{1}(x) = \frac{1}{2}x_{1}^{2} + x_{2}^{2}$,
            \item $V_{1}^{2}(x) = \frac{x_{1}^{2}}{1 - x_{1}x_{2}} + x_{2}^{2}$.
        \end{enumerate}

        %---------------------------------------------------
        \subsubsection{System 2}
        Rozważamy układ:
        \begin{equation}
            \label{eq:zadania_przyklad_2}
            \left\{
            \begin{array}{l}
                 \dot{x}_{1} = x_{2} - x_{1} + x_{1}^{3}\\
                 \dot{x}_{2} = -x_{1}
            \end{array}
            \right\}.
        \end{equation}
        Dla tego układu rozpatrujemy dwóch kandydatów na funkcję Lapunowa:
        \begin{enumerate}[label=(\alph*)]
            \item $V_{2}^{1}(x) = \frac{1}{2} \left( x_{1}^{2} + x_{2}^{2} \right)$,
            \item $V_{2}^{2}(x) = \frac{1}{2} \left( 2x_{1}^{2} + 2x_{1}x_{2} + 3x_{2}^{2} \right)$.
        \end{enumerate}

        %---------------------------------------------------
        \subsubsection{System 3}
        Rozważamy układ (zauważmy, że nie rozważamy $x\in \mathbb{R}^{2}$, ale $x_{1} > -1$):
        \begin{equation}
            \label{eq:zadania_przyklad_3}
            \left\{
            \begin{array}{l}
                 \dot{x}_{1} = x_{2}\\
                 \dot{x}_{2} = -\frac{x_{2}}{1 + x_{1}} - \frac{x_{1}}{1 + x_{1}}
            \end{array}
            \right\}.
        \end{equation}
        Dla tego układu rozpatrujemy dwóch kandydatów na funkcję Lapunowa:
        \begin{enumerate}[label=(\alph*)]
            \item $V_{3}^{1}(x) = \frac{1}{2}x_{2}^{2} + \int_{0}^{x_{1}}{\frac{z}{1+z}dz}$,
            \item $V_{3}^{2}(x) = x_{1}^{2} + x_{1}x_{2} + x_{2}^{2}$.
        \end{enumerate}

        \subsubsection{System 4}
        Rozważamy układ (szukamy cyklu granicznego):
        \begin{equation}
            \label{eq:zadania_przyklad_4}
            \left\{
            \begin{array}{l}
                 \dot{x}_{1} = x_{2} - x_{1}^7 \left( x_{1}^{4} + 2x_{2}^{2} - 10 \right)\\
                 \dot{x}_{2} = -x_{1}^{3} - 3x_{2}^{5} \left( x_{1}^{4} + 2x_{2}^{2} - 10 \right)
            \end{array}
            \right\}.
        \end{equation}
        Dla tego układu rozpatrujemy następującego kandydata na funkcję Lapunowa:
        \begin{enumerate}[label=(\alph*)]
            \item $V_{4}^{1}(x) = \left( x_{1}^{4} + 2x_{2}^{2} - 10 \right)^{2}$.
        \end{enumerate}
\end{document}
